Diffusion models under-produce high spatial frequencies, and a growing family
of training objectives sets out to correct this by weighting wavelet or Fourier
sub-bands of the loss. We audit one such objective over five seeds and find
that although it does exactly what it claims---closing \correction~dB of a
\deficitHigh~dB high-frequency deficit on CIFAR-10---it is no better than plain
MSE on FID, Inception Score or KID, and slightly worse on all three.

We then quantify what that correction should have been worth. Applying the
identical correction to the baseline's own samples as a Fourier post-process
improves FID by $\boostValue \pm \boostCI$, so the trained objective ends
\shortfall~FID short of what its own mechanism predicts: it buys a real
spectral improvement and pays for it several times over elsewhere. A twenty-line
post-process on the baseline outperforms the trained objective by that same
margin.

We further show that the published weighting conflates frequency balance with a
Min-SNR-style timestep reweighting ($|r| = 0.82$); isolating the frequency
component while holding the gradient budget fixed makes results significantly
worse ($+1.34$~FID, $p = 0.003$). Finally we report two measurement hazards that
reversed conclusions during this study: a hand-written KID estimator that
inverted the sign of the only positive result, and the observation that
calibrating metric sensitivity on real images measures it at a minimum where it
is quadratically insensitive, understating operating-point sensitivity roughly
thirtyfold.
